\documentclass[aspectratio=169,11pt]{beamer}
% Shared Beamer style for practice contest solution walkthroughs.
\usepackage[utf8]{inputenc}
\usepackage[T1]{fontenc}
\usepackage{lmodern}
\usepackage{amsmath,amssymb}
\usepackage{xcolor}
\usepackage{hyperref}
\usepackage{listings}

\definecolor{CPNavy}{HTML}{172033}
\definecolor{CPBlue}{HTML}{2563EB}
\definecolor{CPGreen}{HTML}{16A34A}
\definecolor{CPOrange}{HTML}{F97316}
\definecolor{CPGray}{HTML}{64748B}
\definecolor{CPLight}{HTML}{F8FAFC}
\definecolor{CPCodeBg}{HTML}{F1F5F9}
\definecolor{CPCodeRule}{HTML}{CBD5E1}

\usetheme{default}
\usefonttheme{professionalfonts}
\setbeamertemplate{navigation symbols}{}
\setbeamersize{text margin left=0.65cm,text margin right=0.65cm}
\setbeamertemplate{itemize item}{\color{CPBlue}$\blacktriangleright$}
\setbeamertemplate{itemize subitem}{\color{CPGray}$\triangleright$}

\setbeamercolor{normal text}{fg=CPNavy,bg=white}
\setbeamercolor{frametitle}{fg=white,bg=CPNavy}
\setbeamercolor{title}{fg=white,bg=CPNavy}
\setbeamercolor{subtitle}{fg=CPLight,bg=CPNavy}
\hypersetup{colorlinks=true,urlcolor=CPBlue}

\setbeamertemplate{frametitle}{%
  \nointerlineskip%
  \begin{beamercolorbox}[wd=\paperwidth,ht=0.88cm,dp=0.22cm,leftskip=0.55cm,rightskip=0.35cm]{frametitle}
    \usebeamerfont{frametitle}\insertframetitle\hfill\small\insertframenumber/\inserttotalframenumber
  \end{beamercolorbox}%
}

\setbeamertemplate{footline}{%
  \leavevmode%
  \hbox{%
    \begin{beamercolorbox}[wd=.58\paperwidth,ht=2.4ex,dp=1ex,leftskip=.5cm]{author in head/foot}%
      \color{CPGray}\insertshorttitle
    \end{beamercolorbox}%
    \begin{beamercolorbox}[wd=.42\paperwidth,ht=2.4ex,dp=1ex,rightskip=.5cm plus1fil]{date in head/foot}%
      \hfill\color{CPGray}\insertshortauthor
    \end{beamercolorbox}%
  }%
}

\setbeamertemplate{title page}{%
  \vspace*{1.25cm}
  \begin{beamercolorbox}[wd=\paperwidth,sep=0.65cm,leftskip=0.15cm,rightskip=0.15cm]{title}
    {\color{white}\Huge\bfseries\inserttitle\par}
    \vspace{0.35cm}
    {\color{CPLight}\Large\insertsubtitle\par}
    \vspace{0.6cm}
    {\color{CPLight}\large\insertauthor\par}
  \end{beamercolorbox}
  \vfill
}

\newcommand{\sectiontag}[1]{\textbf{\color{CPBlue}#1}}
\newenvironment{tightitemize}{\begin{itemize}\small\setlength{\itemsep}{0.18cm}}{\end{itemize}}
\lstdefinestyle{cpcode}{
  language=Python,
  basicstyle=\ttfamily\tiny,
  keywordstyle=\color{CPBlue}\bfseries,
  commentstyle=\color{CPGray}\itshape,
  stringstyle=\color{CPGreen},
  backgroundcolor=\color{CPCodeBg},
  frame=single,
  framerule=0.25pt,
  rulecolor=\color{CPCodeRule},
  showstringspaces=false,
  columns=fullflexible,
  keepspaces=true,
  breaklines=true,
  breakatwhitespace=false,
  tabsize=4,
  xleftmargin=0.08cm,
  xrightmargin=0.08cm,
  aboveskip=0.08cm,
  belowskip=0.08cm
}


\title[classrooms]{Classrooms}
\subtitle{Day 3 solution walkthrough}
\author[Solution Deck]{Competitive Programming Bootcamp}
\date{}

\begin{document}

\begin{frame}[plain]
  \titlepage
\end{frame}

% BEGIN GENERATED STATEMENT INTERPRETATION
\begin{frame}{Interpret the Word Problem}
\small
\sectiontag{Statement excerpts:}
\begin{tightitemize}
  \item \textit{``classrooms on campus''}
  \item \textit{``as many activities as possible''}
\end{tightitemize}

\vspace{0.18cm}
\sectiontag{Programming interpretation:}
\begin{tightitemize}
  \item This is interval scheduling with up to k simultaneous rooms.
  \item When more than k activities overlap, discard the one that ends latest.
  \item That keeps earlier finishing activities and leaves more room for the future.
\end{tightitemize}
\end{frame}
% END GENERATED STATEMENT INTERPRETATION







\begin{frame}{Problem Model}
\small
\sectiontag{Textbook connection:} CP4 3.4 a. Classical Greedy

\vspace{0.25cm}
\sectiontag{Core technique:} Greedy interval scheduling with k rooms

\vspace{0.25cm}
\begin{tightitemize}
  \item Activities have start and end times.
  \item At most k activities can be active at the same time.
  \item Activities ending exactly at another start still overlap in this problem.
\end{tightitemize}
\end{frame}

\begin{frame}{How to Recognize the Approach}
\begin{tightitemize}
  \item This is interval scheduling with multiple available rooms.
  \item When too many intervals overlap, keep the ones that end earliest.
  \item Priority queues let us remove finished intervals and reject the worst active interval.
\end{tightitemize}
\end{frame}

\begin{frame}{Algorithm}
\begin{tightitemize}
  \item Sort activities by start time.
  \item Use a min-heap to expire active activities with end < current start.
  \item Tentatively accept the current activity.
  \item Use a max-heap of ending times to remove the active activity that ends latest if active > k.
  \item Count accepted activities.
\end{tightitemize}
\end{frame}

\begin{frame}{Walkthrough of the Key State}
\begin{tightitemize}
  \item status marks whether an activity is active, removed, or finished.
  \item Lazy deletion handles heap entries whose status has changed.
  \item Removing the latest-ending active interval leaves the most room for future starts.
\end{tightitemize}
\end{frame}

\begin{frame}{Why the Algorithm Is Correct}
\begin{tightitemize}
  \item At every start time, the accepted active set is feasible after any needed removal.
  \item If one active interval must be removed, dropping the one with latest end cannot reduce future availability compared with dropping an earlier-ending one.
  \item By exchange, the greedy active set can be transformed into an optimal active set after each step.
\end{tightitemize}
\end{frame}

\begin{frame}{Complexity and Implementation Notes}
\small
\sectiontag{Complexity:} O(N log N) time and O(N) memory.

\vspace{0.3cm}
\begin{tightitemize}
  \item Use end < start, not end <= start, because equal endpoints overlap here.
  \item Python max-heap behavior is simulated with negative end times.
\end{tightitemize}
\end{frame}



% BEGIN GENERATED CODE WALKTHROUGH
\begin{frame}[fragile]{Code: Read and Sort Intervals}
\scriptsize
\sectiontag{Source lines:} \texttt{classrooms.py:40--56}

\vspace{0.08cm}
\begin{columns}[T,totalwidth=\textwidth]
\begin{column}{0.58\textwidth}
\begin{lstlisting}[style=cpcode]
import sys
import heapq

def main():
    input = sys.stdin.buffer.readline

    n, k = map(int, input().split())

    activities = []

    for i in range(n):
        start, end = map(int, input().split())
        activities.append((start, end, i))

    activities.sort()
\end{lstlisting}
\end{column}
\begin{column}{0.38\textwidth}
\sectiontag{What this chunk does}
\begin{tightitemize}
  \item Each activity stores start, end, and original index.
  \item Sorting by start time lets the sweep process events chronologically.
  \item The original index is used to keep heap entries tied to status values.
\end{tightitemize}
\end{column}
\end{columns}
\end{frame}

\begin{frame}[fragile]{Code: Track Active Choices}
\scriptsize
\sectiontag{Source lines:} \texttt{classrooms.py:58--72}

\vspace{0.08cm}
\begin{columns}[T,totalwidth=\textwidth]
\begin{column}{0.58\textwidth}
\begin{lstlisting}[style=cpcode]
min_heap = []

max_heap = []

status = [0] * n

active = 0
accepted = 0
\end{lstlisting}
\end{column}
\begin{column}{0.38\textwidth}
\sectiontag{What this chunk does}
\begin{tightitemize}
  \item The min-heap identifies accepted activities that have already ended.
  \item The max-heap identifies the currently active activity with latest ending time.
  \item status handles stale heap entries after an activity has been removed or finished.
\end{tightitemize}
\end{column}
\end{columns}
\end{frame}

\begin{frame}[fragile]{Code: Free Finished Rooms}
\scriptsize
\sectiontag{Source lines:} \texttt{classrooms.py:74--82}

\vspace{0.08cm}
\begin{columns}[T,totalwidth=\textwidth]
\begin{column}{0.58\textwidth}
\begin{lstlisting}[style=cpcode]
for start, end, index in activities:

    while min_heap and min_heap[0][0] < start:
        old_end, old_index = heapq.heappop(min_heap)

        if status[old_index] == 1:
            status[old_index] = 2
            active -= 1
\end{lstlisting}
\end{column}
\begin{column}{0.38\textwidth}
\sectiontag{What this chunk does}
\begin{tightitemize}
  \item Before accepting a new activity, finished activities stop occupying rooms.
  \item The strict check old\_end < start matches the statement that equal endpoints still overlap.
  \item Only status 1 activities affect the active-room count.
\end{tightitemize}
\end{column}
\end{columns}
\end{frame}

\begin{frame}[fragile]{Code: Accept or Drop Latest}
\scriptsize
\sectiontag{Source lines:} \texttt{classrooms.py:84--104}

\vspace{0.08cm}
\begin{columns}[T,totalwidth=\textwidth]
\begin{column}{0.58\textwidth}
\begin{lstlisting}[style=cpcode]
status[index] = 1

heapq.heappush(min_heap, (end, index))
heapq.heappush(max_heap, (-end, index))

active += 1
accepted += 1

if active > k:

    while max_heap:
        neg_end, remove_index = heapq.heappop(max_heap)

        if status[remove_index] == 1:
            status[remove_index] = 0
            active -= 1
            accepted -= 1
            break
\end{lstlisting}
\end{column}
\begin{column}{0.38\textwidth}
\sectiontag{What this chunk does}
\begin{tightitemize}
  \item The new activity is tentatively accepted first.
  \item If active exceeds k, the latest-ending active activity is removed.
  \item Removing the latest finisher is the greedy choice that preserves the most flexibility.
\end{tightitemize}
\end{column}
\end{columns}
\end{frame}

\begin{frame}[fragile]{Code: Print Accepted Count}
\scriptsize
\sectiontag{Source lines:} \texttt{classrooms.py:106--109}

\vspace{0.08cm}
\begin{columns}[T,totalwidth=\textwidth]
\begin{column}{0.58\textwidth}
\begin{lstlisting}[style=cpcode]
    print(accepted)

main()
\end{lstlisting}
\end{column}
\begin{column}{0.38\textwidth}
\sectiontag{What this chunk does}
\begin{tightitemize}
  \item accepted has been updated throughout the sweep.
  \item Only the maximum count is required, not the actual classroom assignment.
\end{tightitemize}
\end{column}
\end{columns}
\end{frame}
% END GENERATED CODE WALKTHROUGH

\begin{frame}{Source}
\small
\sectiontag{Solution file:} \texttt{practice\_contests/day\_3/classrooms.py}

\vspace{0.3cm}
\sectiontag{Problem:} \url{https://open.kattis.com/problems/classrooms}

\vspace{0.45cm}
Use this deck with the source code open beside it. The slides explain the reasoning; the code shows the exact input handling and output formatting.
\end{frame}

\end{document}
